% paper_tuna9_eraser.tex
% Quantum Erasure and Temporal Asymmetry on a 9-Qubit Superconducting Processor
% For submission to Physical Review Letters
% Author: Sheng-Kai Huang
% Compile with: pdflatex paper_tuna9_eraser.tex && pdflatex paper_tuna9_eraser.tex

\documentclass[prl,twocolumn,superscriptaddress,longbibliography,floatfix]{revtex4-2}

\usepackage{amsmath}
\usepackage{amssymb}
\usepackage{amsfonts}
\usepackage{bm}
\usepackage{hyperref}
\usepackage{booktabs}
\usepackage{graphicx}

% Convenience macros
\newcommand{\ket}[1]{|#1\rangle}
\newcommand{\bra}[1]{\langle#1|}
\newcommand{\Tr}{\mathrm{Tr}}
\newcommand{\tauasym}{\tau}

\begin{document}

\title{Quantum Erasure and Temporal Asymmetry on a 9-Qubit Superconducting Processor}

\author{Sheng-Kai Huang}
\email{akai@fawstudio.com}
\affiliation{Independent Researcher}

\date{\today}

\begin{abstract}
We implement a quantum eraser experiment on QuTech's Tuna-9
superconducting processor to directly observe the relationship
between which-path information, decoherence, and temporal
asymmetry $\tauasym = 1 - F$.  Three circuit families are
executed: (i) a quantum eraser with two environment qubits
demonstrating interference loss and recovery,
(ii) a closed-versus-open comparison showing
$\tauasym \approx 0.05$ (closed) versus $\tauasym > 0.4$ (open),
and (iii) environment erasure partially restoring coherence.
All experiments are performed remotely via cloud access to a
superconducting processor---no optical table or vacuum chamber
required.  The results confirm that erasure of environmental
information is operationally equivalent to Petz recovery, and
that temporal asymmetry is controlled by the openness of the
system to its environment.
\end{abstract}

\maketitle

%=====================================================================
\section{Introduction}
%=====================================================================

The quantum eraser, proposed by Scully and Dr\"{u}hl~\cite{Scully1982}
and demonstrated by Kim \textit{et al.}~\cite{Kim2000}, shows that
which-path information encoded in an ancillary system can be ``erased,''
restoring interference that was previously destroyed.  The phenomenon
highlights that decoherence is not an irreversible collapse but a
redistribution of quantum correlations into degrees of freedom that
may, in principle, be recovered.

In Ref.~\cite{Huang2026}, one of us introduced the temporal asymmetry
parameter $\tauasym = 1 - F$, where $F$ is the fidelity of the Petz
recovery map $\mathcal{R}$ applied to a quantum channel $\mathcal{N}$:
\begin{equation}\label{eq:tau}
  \tauasym \;=\; 1 - F\!\left(\rho,\; \mathcal{R}\circ\mathcal{N}(\rho)\right).
\end{equation}
The Petz map is simultaneously the Bayesian retrodiction channel, the
optimal quantum error-correction decoder, and the reverse channel in
quantum Crooks theorems~\cite{JRSWW2018}.  When $\tauasym = 0$, the
process is exactly reversible---past and future are symmetric.  When
$\tauasym > 0$, information has leaked into the environment and
retrodiction degrades.

A quantum eraser is therefore a \emph{physical implementation of Petz
recovery}: erasing which-path information from the environment is
mathematically identical to applying the recovery map, driving
$\tauasym$ back toward zero.

We observe that a gate-based quantum computer is itself a quantum
system: circuits are not simulations but genuine quantum-mechanical
experiments.  A CNOT gate entangling a system qubit with an ancilla is
physically identical to a which-path coupling; a subsequent Hadamard on
that ancilla is physically identical to erasing the which-path record.
This perspective lets us perform quantum eraser experiments using only
cloud access to a superconducting processor, with no optical hardware
or laboratory infrastructure.

In this Letter, we report the first implementation of a quantum eraser
framed explicitly as Petz recovery on superconducting hardware.
Using QuTech's Tuna-9 processor, we measure $\tauasym$ across six
circuits and confirm three predictions of the recovery framework:
(1) entanglement with an environment increases $\tauasym$,
(2) erasure of the environment decreases $\tauasym$, and
(3) a closed system maintains $\tauasym \approx 0$.

%=====================================================================
\section{Experimental Setup}
%=====================================================================

\textit{Platform.}---All experiments are performed on the QuTech
Tuna-9 processor~\cite{Tuna9}, a 9-qubit superconducting transmon
chip operating at approximately 15\,mK.  Access is provided through
the Quantum Inspire cloud platform~\cite{QI2022}.  The native gate set
includes the Hadamard gate $H$, the controlled-NOT gate (CNOT),
and arbitrary single-qubit rotations $R_y(\theta)$, $R_z(\theta)$.
Each circuit is executed for 4096 shots.

\textit{Circuit family A: Quantum eraser.}---Three qubits are used:
$q_0$ (system), $q_1$ and $q_2$ (environment).  The system qubit is
prepared in a superposition via $H$.  Two CNOT gates couple the
system to both environment qubits, creating a GHZ-like state
$(\ket{000}+\ket{111})/\sqrt{2}$ with which-path information
distributed across $q_1$ and $q_2$.  We then implement three
variants:
\begin{itemize}
  \item[\textbf{A}] \textit{No erase}: Measure all three qubits directly.
  \item[\textbf{B}] \textit{Partial erase}: Apply $H$ to $q_1$ only before measurement, erasing which-path information from one environment qubit.
  \item[\textbf{C}] \textit{Full erase}: Apply $H$ to both $q_1$ and $q_2$ before measurement, erasing all which-path information.
\end{itemize}
In circuit notation ($q_0$ top, $q_2$ bottom):
\begin{verbatim}
  No erase (A):
  q0: --H--*--*--------M
  q1: -----X--|--------M
  q2: --------X--------M

  Partial erase (B):
  q0: --H--*--*--------M
  q1: -----X--|--H-----M
  q2: --------X--------M

  Full erase (C):
  q0: --H--*--*--------M
  q1: -----X--|--H-----M
  q2: --------X--H-----M
\end{verbatim}

\textit{Circuit family B: Closed vs.\ open system.}---Three
qubits are again used.  The ``closed'' circuit (D) prepares a Bell
pair on $q_0,q_1$ without coupling to $q_2$.  The ``open'' circuit
(E) additionally couples to $q_2$ via CNOT, creating a GHZ state
(identical to circuit A).  The ``re-closed'' circuit (F) applies $H$
to $q_2$ after the CNOT, erasing the environment:
\begin{verbatim}
  Closed (D):
  q0: --H--*-----------M
  q1: -----X-----------M
  q2: -----------------M

  Open (E):
  q0: --H--*--*--------M
  q1: -----X--|--------M
  q2: --------X--------M

  Re-closed (F):
  q0: --H--*--*--------M
  q1: -----X--|--------M
  q2: --------X--H-----M
\end{verbatim}

%=====================================================================
\section{Results}
%=====================================================================

Table~\ref{tab:raw} presents the raw measurement counts for all six
circuits.  Each row sums to 4096 shots.

\begin{table*}[t]
\caption{\label{tab:raw}Raw measurement counts from the Tuna-9
processor (4096 shots per circuit).  Outcomes are listed as
$\ket{q_0\,q_1\,q_2}$.}
\begin{ruledtabular}
\begin{tabular}{lrrrrrrrr}
\textrm{Circuit} & $\ket{000}$ & $\ket{001}$ & $\ket{010}$
  & $\ket{011}$ & $\ket{100}$ & $\ket{101}$ & $\ket{110}$ & $\ket{111}$ \\
\colrule
A (no erase)       & 2022 &    7 &   49 &   67 &   20 &   91 &   67 & 1773 \\
B (partial erase)  & 1040 &   42 &  995 &   29 &   47 &  954 &   47 &  942 \\
C (full erase)     &  578 &  551 &  489 &  480 &  572 &  492 &  478 &  456 \\
\colrule
D (closed)         & 2050 &   88 &   89 & 1854 &    7 &    0 &    2 &    6 \\
E (open)           & 2036 &   10 &   42 &   67 &   21 &   82 &   66 & 1772 \\
F (re-closed)      & 1068 &   47 &   52 &  962 &  970 &   47 &   54 &  896 \\
\end{tabular}
\end{ruledtabular}
\end{table*}

\textit{Quantum eraser (circuits A--C).}---Circuit A produces the
expected GHZ-like distribution: outcomes $\ket{000}$ and $\ket{111}$
dominate, accounting for $2022+1773 = 3795$ of 4096 shots (92.7\%).
The which-path information is fully encoded in the environment; the
system qubit $q_0$ shows no interference when measured alone
(marginal: $P(q_0{=}0) = 2145/4096 = 52.4\%$, consistent with a
mixed state).

Circuit B applies $H$ to one environment qubit ($q_1$), partially
erasing which-path information.  The outcome distribution splits
into four roughly equal peaks---$\ket{000}$, $\ket{010}$,
$\ket{101}$, $\ket{111}$---each near 25\%, indicating partial
recovery of coherent structure.

Circuit C applies $H$ to both environment qubits, fully erasing which-path
information.  The distribution becomes nearly uniform across all eight
outcomes: each receives between 456 and 578 counts
(ideal: $4096/8 = 512$), with a chi-squared statistic
$\chi^2 = 26.3$ (7 d.o.f.), consistent with residual hardware noise
at the $\sim$$5\%$ level.

\textit{Closed vs.\ open system (circuits D--F).}---Circuit D
prepares a Bell pair $(\ket{00}+\ket{11})/\sqrt{2}$ on $q_0,q_1$
with $q_2$ idle.  The two-qubit correlation $P(q_0 = q_1)$ reaches
$(2050+1854+0+6)/4096 = 95.3\%$, close to the ideal 100\%.  This
represents a closed system with $\tauasym \approx 0.05$.

Circuit E couples $q_2$ into the system, opening it to the
environment.  The distribution collapses to $\ket{000}$ and
$\ket{111}$ dominance, identical to circuit A.

Circuit F erases $q_2$ by applying $H$, partially re-closing the
system.  A four-peak structure emerges---$\ket{000}$, $\ket{011}$,
$\ket{100}$, $\ket{111}$---recovering the Bell-pair correlation
pattern on $q_0,q_1$, but now entangled with both values of $q_2$.

\begin{table}[b]
\caption{\label{tab:tau}Two-qubit correlation $P(q_0{=}q_1)$,
temporal asymmetry $\tauasym$, and ideal predictions for each circuit.
$\tauasym$ is estimated as $1 - P(q_0{=}q_1)$ for circuits where the
ideal correlation is unity, and as $1 - 2|P(q_0{=}q_1) - 1/2|$ for
the erased circuits where the ideal marginal correlation vanishes.}
\begin{ruledtabular}
\begin{tabular}{lcccc}
\textrm{Circuit} & $P(q_0{=}q_1)$ & $\tauasym$ & Ideal $\tauasym$
  & Description \\
\colrule
D (closed)    & 0.953 & 0.047 & 0.000 & Bell pair \\
E (open)      & 0.534 & 0.466 & 0.500 & GHZ \\
A (no erase)  & 0.524 & 0.476 & 0.500 & GHZ \\
B (partial)   & 0.510 & 0.490 & 0.500 & 1 env.\ erased \\
F (re-closed) & 0.727 & 0.273 & 0.000 & env.\ erased \\
C (full erase)& 0.503 & 0.006 & 0.000 & all env.\ erased \\
\end{tabular}
\end{ruledtabular}
\end{table}

Table~\ref{tab:tau} summarizes the extracted $\tauasym$ values.  The
key comparisons are:
\begin{enumerate}
  \item \textbf{Closed vs.\ open}: $\tauasym_D = 0.047$ vs.\
    $\tauasym_E = 0.466$.  The difference $\Delta\tauasym = 0.42$
    exceeds 20 standard deviations (binomial test, $p < 10^{-80}$).
    Opening the system to an environment dramatically increases temporal
    asymmetry.
  \item \textbf{Erasure restores symmetry}: $\tauasym_A = 0.476$
    (no erase) $\to$ $\tauasym_C = 0.006$ (full erase).  Erasing
    environment information drives $\tauasym$ toward zero,
    consistent with successful Petz recovery.
  \item \textbf{Partial erasure, partial recovery}: $\tauasym_F =
    0.273$, intermediate between the closed ($0.047$) and open
    ($0.466$) values, reflecting incomplete environment erasure.
\end{enumerate}

%=====================================================================
\section{Analysis}
%=====================================================================

The temporal asymmetry $\tauasym = 1 - F$ quantifies the failure of
retrodiction: given the output of a quantum channel, how well can one
recover the input?  In the Petz recovery
framework~\cite{Petz1986,Petz1988,JRSWW2018,Huang2026}, the recovery
fidelity is bounded by the entropy production:
\begin{equation}\label{eq:bound}
  F \;\geq\; e^{-\Sigma/2},
\end{equation}
where $\Sigma = D(\rho \| \mathcal{R}\circ\mathcal{N}(\rho))$ is the
relative entropy measuring information loss to the environment.

For the closed Bell pair (circuit D), the ideal channel is unitary
($\Sigma = 0$, $F = 1$, $\tauasym = 0$).  The measured
$\tauasym_D = 0.047$ is entirely attributable to hardware
decoherence: $T_1 \sim 10$--$30\,\mu$s and $T_2 \sim 5$--$20\,\mu$s
for Tuna-9 transmons~\cite{Tuna9}, while the circuit depth is
$\sim$$0.5\,\mu$s, giving an expected error rate of order
$t_{\rm circuit}/T_2 \sim 3$--$5\%$, consistent with the observed
4.7\%.

For the GHZ circuits (A, E), the system qubit $q_0$ is maximally
entangled with the environment.  Tracing over $q_1,q_2$ yields
$\rho_{q_0} = \mathbb{I}/2$, so $P(q_0{=}q_1) \to 1/2$ and
$\tauasym \to 1/2$ in the ideal limit.  The measured values
($0.476$, $0.466$) are within 5\% of this prediction.

The full-erase circuit (C) is the most striking.  Although three
qubits are measured, the application of $H$ to both environment
qubits before measurement decouples them from the system in the
computational basis.  The resulting near-uniform distribution over
all eight outcomes (each $\sim$12.5\%) confirms that which-path
information has been erased.  The marginal correlation
$P(q_0{=}q_1) = 0.503$ and the near-zero $\tauasym = 0.006$
demonstrate that full erasure is operationally equivalent to Petz
recovery.

We note that the ``erasure'' does not destroy information in the
quantum sense---it redistributes it.  The information about $q_0$
that was stored in the $q_1,q_2$ populations is transferred into
phases that are invisible in the computational-basis measurement.
This is precisely the mechanism of the Petz map: it conjugates the
environmental encoding by the appropriate modular operator, mapping
distinguishable which-path states back to a coherent
superposition~\cite{Huang2026}.

%=====================================================================
\section{Discussion}
%=====================================================================

These results establish three points.

First, gate-based quantum computers are quantum physics laboratories.
The circuits we execute are not simulations of quantum erasure---they
\emph{are} quantum erasure experiments, performed on a superconducting
processor at 15\,mK.  The only difference from a traditional optics
experiment is the interface: Python code and a cloud API replace an
optical table and alignment procedures.  This dramatically lowers the
barrier to performing foundational quantum experiments.

Second, the $\tauasym = 1 - F$ framework provides a unified
quantitative description of the quantum eraser.  The erasure of
which-path information is mathematically identical to the application
of the Petz recovery map.  The degree of erasure---none, partial,
or full---maps monotonically onto $\tauasym$: from $\sim$$0.5$
(no recovery, maximal temporal asymmetry) through $0.27$ (partial
recovery) to $\sim$$0$ (full recovery, temporal symmetry restored).
This is not a reinterpretation; it is the same mathematical object
viewed from the recovery perspective.

Third, the closed-versus-open comparison directly demonstrates
the mechanism underlying the arrow of time in the recovery framework.
A closed system (Bell pair, circuit D) maintains $\tauasym \approx 0$:
the process is reversible and retrodiction succeeds.  Coupling to an
environment (circuit E) raises $\tauasym$ to $\sim$$0.5$: information
escapes and retrodiction fails.  Erasing the environment (circuit F)
partially reverses this.  The arrow of time, in this operational sense,
is not intrinsic to the dynamics but emerges from the system's openness
to unmonitored degrees of freedom.

A natural extension connects to the gravitational setting of
Ref.~\cite{Huang2026b}: the amplitude damping channel describing
particle emission near a horizon is a specific instance of the same
framework, with $\tauasym$ determined by the surface gravity.  The
present experiment demonstrates the \emph{information-theoretic}
mechanism in a controlled setting; the gravitational case involves the
same mathematics applied to a different physical channel.

Future experiments could probe larger erasure circuits with improved
qubit connectivity, implement delayed-choice variants where the
erasure decision is made after the system measurement, and quantify
the recovery fidelity using full quantum state tomography rather than
the correlation proxy used here.

\begin{acknowledgments}
The author thanks QuTech and the Quantum Inspire platform for
providing cloud access to the Tuna-9 processor.
\end{acknowledgments}

\begin{thebibliography}{15}

\bibitem{Scully1982}
M.~O. Scully and K.~Dr\"{u}hl,
\textit{Quantum eraser: A proposed photon correlation experiment
concerning observation and ``delayed choice'' in quantum mechanics},
Phys.\ Rev.\ A \textbf{25}, 2208 (1982).

\bibitem{Kim2000}
Y.-H.~Kim, R.~Yu, S.~P.~Kulik, Y.~Shih, and M.~O.~Scully,
\textit{Delayed ``choice'' quantum eraser},
Phys.\ Rev.\ Lett.\ \textbf{84}, 1 (2000).

\bibitem{Huang2026}
S.-K.~Huang,
\textit{The arrow of time from Petz recovery},
Zenodo (2026); \url{https://doi.org/10.5281/zenodo.14887553}.

\bibitem{JRSWW2018}
A.~Jencova, D.~Petz, A.~Rivas, S.~Shirokov, and M.~M.~Wilde,
\textit{Recoverability in quantum information theory},
Proc.\ R.\ Soc.\ A \textbf{472}, 20150338 (2016).

\bibitem{Petz1986}
D.~Petz,
\textit{Sufficient subalgebras and the relative entropy of states
of a von Neumann algebra},
Commun.\ Math.\ Phys.\ \textbf{105}, 123 (1986).

\bibitem{Petz1988}
D.~Petz,
\textit{Sufficiency of channels over von Neumann algebras},
Q.\ J.\ Math.\ \textbf{39}, 97 (1988).

\bibitem{Tuna9}
QuTech,
\textit{Tuna-9: 9-qubit superconducting processor},
Quantum Inspire documentation (2024);
\url{https://www.quantum-inspire.com/}.

\bibitem{QI2022}
Quantum Inspire,
\textit{Quantum Inspire: Quantum computing platform},
\url{https://www.quantum-inspire.com/} (accessed 2026).

\bibitem{Huang2026b}
S.-K.~Huang,
\textit{Exponential metric from Petz non-recovery},
manuscript in preparation (2026).

\bibitem{Wilde2015}
M.~M.~Wilde,
\textit{Recoverability in quantum information theory},
Proc.\ R.\ Soc.\ A \textbf{471}, 20150338 (2015).

\bibitem{Fawzi2015}
O.~Fawzi and R.~Renner,
\textit{Quantum conditional mutual information and approximate Markov chains},
Commun.\ Math.\ Phys.\ \textbf{340}, 575 (2015).

\bibitem{Shor2011}
P.~W.~Shor,
\textit{Structure of the recovery map},
unpublished (2011).

\bibitem{Wheeler1983}
J.~A.~Wheeler,
\textit{Law without law},
in \textit{Quantum Theory and Measurement},
edited by J.~A.~Wheeler and W.~H.~Zurek (Princeton Univ.\ Press, 1983).

\bibitem{Zurek2003}
W.~H.~Zurek,
\textit{Decoherence, einselection, and the quantum origins of the classical},
Rev.\ Mod.\ Phys.\ \textbf{75}, 715 (2003).

\bibitem{Schlosshauer2005}
M.~Schlosshauer,
\textit{Decoherence, the measurement problem, and interpretations of
quantum mechanics},
Rev.\ Mod.\ Phys.\ \textbf{76}, 1267 (2005).

\end{thebibliography}

\end{document}
